\begin{frame}{kNN과의 차이 --- 그리고 단순함의 가치}
  \begin{columns}[T]
    \begin{column}{0.5\textwidth}
      \kb{kNN}(게으름, lazy)
      \begin{itemize}\small
        \item 예측할 때마다 \textbf{그때그때} 거리 재기
        \item 거리 계산이 \textbf{예측 시점}에 몰려 있음
      \end{itemize}
    \end{column}
    \begin{column}{0.5\textwidth}
      \kb{k-means}
      \begin{itemize}\small
        \item 중심점이 \textbf{수렴할 때까지} 데이터를
          반복적으로 훑음
        \item 거리 계산이 ``\textbf{학습} 단계''에 몰려 있음
          $\Rightarrow$ 오히려 \textbf{선형회귀 같은 파라미터
          학습} 쪽에 더 가깝다
        \item \textbf{라벨이 없다는 것만} 지도학습과 다름
      \end{itemize}
    \end{column}
  \end{columns}
  \vskip0.8em
  \begin{block}{단순함이 곧 약점이 아니다}
    kNN과 k-means 둘 다 모델이 ``무엇을 학습했는지''를 설명해주는
    복잡한 이론이 없지만, \textbf{바로 그 단순함}(순전히 거리
    계산) 덕분에 \textbf{언제, 왜 실패하는지도 수식으로
    정확히} 짚어낼 수 있다.
  \end{block}
\end{frame}
